\chapter{Related Work}
\label{sec:related-work}

Brzozowski's derivatives, Antimirov's partial derivatives, and Sulzmann
and Lu's POSIX and bit-coded constructions~\citep{brzozowski1964,
antimirov1996, sulzmannlu2014, sulzmannlu2012} are this thesis's direct
basis, presented in
Chapters~\ref{sec:derivatives}--\ref{sec:partial-derivative-parser}.
This chapter places it among the adjacent and alternative approaches
it does not build on.

\textbf{Derivatives beyond POSIX parsing.} Brzozowski's derivatives
saw little use for decades before \citet{owens2009} revived them for
practical lexing and matching, and \citet{mightdaraisspiewak2011}
extended the idea from matching to full context-free parsing. Both
informed the renewed interest in derivative-based techniques this
thesis's approach sits within, but neither targets POSIX
disambiguation or constructs a parse tree the way Sulzmann and Lu's
method does.

\textbf{Greedy disambiguation.} POSIX is one of two disambiguation
policies distinguished in Chapter~\ref{sec:matching-vs-parsing}; the
other, Greedy, has its own efficient algorithm in
\citet{frischcardelli2004}, built from automaton construction rather
than derivatives: a different policy, not a competing route to the
same one.

\textbf{Tagged automata for submatching.} Outside the derivative line,
\citet{laurikari2000} introduced tagged NFAs, which attach submatch
bookkeeping to automaton transitions instead of threading it
symbolically through a derivative or bit-code construction.
Kuklewicz's \texttt{regex-tdfa}~\citep{kuklewicz_tdfa}, a production
Haskell implementation of this line, is the dominant non-derivative
architecture for the same submatching problem this thesis's
bit-coded parsers solve symbolically.

\textbf{Production engines.} Outside the academic derivative and
tagged-automaton lines, production engines take a third approach.
RE2~\citep{re2} compiles to a Thompson-style automaton for guaranteed
linear-time matching without constructing a full POSIX parse tree, and
Rust's own \texttt{regex} crate~\citep{rust_regex}, the language this
thesis is implemented in, follows the same lineage. Both prioritize
throughput and backtracking-safety over the parse-tree precision this
thesis targets, approaching the catastrophic-backtracking motivation
of Section~\ref{sec:motivation} from the automaton side rather than
the derivative side.

\textbf{Recent adjacent work.} Two more recent threads sit near this
thesis's approach without overlapping it. \citet{mamouras2024}
analyze a pattern statically for whether its POSIX and Greedy results
can ever diverge, complementing the empirical POSIX-versus-Greedy
comparison in Chapter~\ref{sec:evaluation}.
\citet{varatalu2023} extend derivative-based matching to intersection,
complement, and lookaround, beyond the core language this thesis
implements. Neither targets Rust or a bit-coded construction, leaving
that systems question open before this work.
